\hypertarget{appendix-2e245a6fa8d3}{}\label{referencestart:appendix-2e245a6fa8d3}
\section{Short-circuit evaluation}\label{appendix-2e245a6fa8d3-short-circuit-evaluation}

\hypertarget{appendix-2e245a6fa8d3-short-circuit-evaluation}{}

\textbf{Curriculum depth:} Introduced · \textbf{First seen:} C001 · \textbf{Latest development:} C001

\begin{allmightroledefinition}

\subsection{Definition and mechanism}\label{appendix-2e245a6fa8d3-definition-and-mechanism}

\hypertarget{appendix-2e245a6fa8d3-definition-and-mechanism}{}

\BookInlineCode{\&\&} may skip its right operand after false; \BookInlineCode{\textbar\textbar{}} may skip it after true, so evaluation safety is part of semantics.

\end{allmightroledefinition}

Built-in logical AND skips its right operand after false; logical OR skips it after true. \BookInlineCode{ptr\ \&\&\ ptr-\textgreater ready()} guards the dereference. A bitwise \BookInlineCode{\&} evaluates both operands. Operator spelling therefore determines both result meaning and whether a dangerous expression is evaluated.

\subsection{Prerequisites}\label{appendix-2e245a6fa8d3-prerequisites}

\hypertarget{appendix-2e245a6fa8d3-prerequisites}{}

\begin{itemize}
\tightlist
\item
  \hyperref[appendix-fe6b457d68d8]{Logical versus bitwise operators} --- Built-in \BookInlineCode{\&\&}, \BookInlineCode{\textbar\textbar{}}, and \BookInlineCode{!} convert operands to truth conditions.
\end{itemize}

\subsection{Used by}\label{appendix-2e245a6fa8d3-used-by}

\hypertarget{appendix-2e245a6fa8d3-used-by}{}

\begin{itemize}
\tightlist
\item
  \href{../01-Bits-Programs-and-Reliable-Systems-P02/main.pdf\#nameddest=appendix-652d81568de9}{Checked binary parsing (Bits, Programs, and Reliable Systems, Part 2, Reference notes: Checked binary parsing)} --- uses this mechanism to define its representation, error or validation contract.
\end{itemize}

Connect later chunks here when they use this concept.

\subsection{Recall and next depth}\label{appendix-2e245a6fa8d3-recall-and-next-depth}

\hypertarget{appendix-2e245a6fa8d3-recall-and-next-depth}{}

Explain short-circuit evaluation without the label, work the example above, then state the assumption whose removal would change the result.

\textbf{Pending:} sequencing, side effects, and optimizer control flow.

\subsection{Source and navigation}\label{appendix-2e245a6fa8d3-source-and-navigation}

\hypertarget{appendix-2e245a6fa8d3-source-and-navigation}{}

\hyperref[01-foundations-i-ch001-from-a-signal-to-an-interpretation]{C001 teaching chapter} · \href{../01-Bits-Programs-and-Reliable-Systems-P03/main.pdf\#nameddest=appendix-fd090d982a15}{Part definitions and mechanisms (Bits, Programs, and Reliable Systems, Part 3, Reference notes: Foundations)} · \href{../07-Vision-XR-and-Playable-Worlds-P01/main.pdf\#nameddest=appendix-4ad47bdd269a}{Dependency Map (Vision, XR, and Playable Worlds, Part 1, Reference notes: Dependency Map)}
